\documentclass[12pt,oneside]{article}

\usepackage[margin=1in]{geometry}

\usepackage{setspace}

\usepackage{sectsty}

\usepackage{enumitem}

\usepackage[hidelinks]{hyperref}

\usepackage{times}

\usepackage{footmisc}  % For footnotes

\onehalfspacing

\setlist[itemize]{leftmargin=0.5in}

\sectionfont{\normalsize\bfseries}

\begin{document}

\begin{center}

{\textbf{\large COMPREHENSIVE CRIMINAL REPORT}} \\

{\textbf{E'LISE CHARD}} \\

{\textbf{(YWCA MISSOULA OPERATIONS MANAGER)}} \\

Institutional Coordinator \& Montana Phase Escalator \\

2017--2025

\end{center}

\vspace{0.3in}

\section*{PERPETRATOR PROFILE}

\begin{itemize}

\item \textbf{Name:} E'Lise Chard

\item \textbf{Position:} Operations Manager, YWCA Missoula

\item \textbf{Primary Jurisdiction:} Montana (Missoula County)

\item \textbf{Role:} Institutional coordinator, Montana phase escalation, proxy harassment

\end{itemize}

\section*{RELATIONSHIP HISTORY}

\begin{itemize}

\item Minimal contact with Nuno (2001--2018)

\item Sister to Danielle Chard (primary perpetrator Phase 1)

\item 2001: Told Nuno she had disowned Danielle at age 18

\item Called Danielle ``pathological liar and abusive to men''

\item 2017--2018: Activated Montana campaign when Danielle's credibility exhausted

\end{itemize}

\section*{DOCUMENTED CRIMINAL ACTIONS}

\subsection*{June 12, 2018: Perjurious Protection Order}\footnote{See Exhibit 10, Section 1.4-1.5: ELise-OP-Filing-Page1.jpeg and ELise-OP-Filing-Page2.jpeg (sworn petition pages); Section 2.8: Nuno-YWCA-Complaint.jpg (protected First Amendment petition filed June 3, 2018).}

\textbf{Timeline:} Filed 9 days after Nuno's protected YWCA complaint (\$400 donation)

\textbf{First Amendment Retaliation Pattern:} Temporal proximity establishes causation

\textbf{Written Claims (Petition):}\footnote{Written petition filed under penalty of perjury June 12, 2018, Missoula County. See Exhibit 10, Section 1.4-1.5.}

\begin{itemize}

\item ``Ongoing threats and harassment to her children''

\item ``Extensive criminal drug history known to all police and judges''

\item ``Multiple threatening communications''

\end{itemize}

\textbf{Sworn Testimony (Judge Warren, asked TWICE):}\footnote{Transcript and court findings available in Montana case records. Judge Warren explicitly found Nuno's YWCA complaint constituted protected First Amendment speech.}

\begin{itemize}

\item ``No contact beyond three 2017 emails''

\item ``No threats''

\item Reality: Zero drug history; one coerced plea she tried to withdraw

\end{itemize}

\textbf{Direct Perjury:} Written petition contradicts sworn testimony

\textbf{Federal Crime:} 18 U.S.C. § 1341 (mail fraud via court documents)

\textbf{State Crime:} MCA 45-7-202 (perjury), MCA 45-8-212 (criminal defamation)

\subsection*{July--September 2018: Institutional Conspiracy}

\textbf{Coordination:}

\begin{itemize}

\item Home raid and electronics seizure based on charges

\item Detective Connie Brueckner (YWCA BOARD MEMBER)\footnote{See Exhibit 10, Section 6.1: ywca-missoula-consolidated-financial-statements-years-2023-2024.pdf (ProPublica nonprofit records confirm Brueckner's YWCA board membership). This undisclosed conflict of interest pervades the entire investigation and constitutes institutional capture.} leads investigation

\item Conflict of interest UNDISCLOSED to defense counsel

\item E'Lise worked directly with Brueckner (no court disclosure)

\item Officer Ethan Smith recruited by E'Lise (later removed for ``obvious impropriety'')

\end{itemize}

\textbf{Judicial Outcome:} ALL CHARGES DISMISSED WITH PREJUDICE

\textbf{Judge Warren Finding:} Nuno's YWCA complaint was protected First Amendment speech

\textbf{Federal Crime:} 18 U.S.C. § 242 (deprivation of rights under color of law)

\textbf{State Crime:} MCA 45-4-102 (criminal conspiracy)

\subsection*{August 14, 2018: Suppression of Exculpatory Evidence}\footnote{\textbf{CRITICAL EVIDENCE.} See Exhibit 10, Section 2.7: Nuno-PS-Witness-Statement-YWCA-Email-Complaint.pdf. This August 2018 character reference describes Elvis sheltering a pregnant homeless woman (Witness PS) at no cost in 2016, demonstrating his generous, non-violent character---the exact opposite of E'Lise's false narrative and the exact mission YWCA claims to perform. YWCA and Detective Brueckner suppressed this exculpatory evidence, constituting a Brady violation.}

\textbf{Witness PS Character Reference received by YWCA:}

\begin{quote}

``In approximately 2016, when I was pregnant and homeless, Elvis invited me into his home at no cost. He expected nothing in return. He was never physically or verbally abusive. At times of conflict, he handled himself very well, working through issues instead of making things worse. I hate what your employee E'Lise Chard has been doing to him for no reason.''

\end{quote}\footnote{Email dated August 14, 2018, sent to YWCA Missoula. Never disclosed to Elvis, his defense counsel, or the court despite active criminal proceedings. This constitutes institutional suppression of exculpatory evidence and Brady violation.}

\textbf{YWCA Response:} Complete silence and suppression

\textbf{E'Lise's Knowledge:} As Operations Manager, received/reviewed reference

\textbf{Legal Significance:} Brady violation facilitated by institutional suppression

\textbf{Federal Crime:} 18 U.S.C. § 1503 (obstruction of justice)

\textbf{Civil Liability:} 42 U.S.C. § 1983 (deliberate indifference to civil rights)

\subsection*{September 2018: False Police Reports}\footnote{Subpoenaed phone records showed no calls from Elvis to E'Lise during alleged timeframe. Fabrication documented in Montana prosecution discovery materials.}

\begin{itemize}

\item Allegation: ``Phantom phone calls'' to her

\item Subpoenaed Records: NO SUCH CALLS EXISTED

\item E'Lise fabricated allegations to support escalation

\item Detective Brueckner used fabrication to justify continued prosecution

\item State Crime: MCA 45-5-204 (false reporting)

\end{itemize}

\subsection*{2020--2024: Tyleen Root Recruitment}\footnote{See Exhibit 10, Section 5.1-5.19: Tyleen-Root-Harassment series (19 files documenting defamatory social media posts, business sabotage, and proxy harassment campaign).}

\textbf{Pattern:}

\begin{itemize}

\item E'Lise recruited YWCA client Tyleen Root

\item Directed Tyleen to conduct business sabotage campaign

\item Private DMs → public posts → targeted follower contact

\item False claims: ``Stalker and woman beater''

\item Business Consequence: \$60,000/year contract loss

\end{itemize}

\textbf{Federal Crime:} 18 U.S.C. § 1962 (RICO continuing violation)

\textbf{State Crime:} MCA 45-4-102 (criminal solicitation), MCA 45-8-212 (defamation)

\subsection*{2025: Ongoing Retaliation}

\begin{itemize}

\item Continues harassment through proxies

\item Institutional protection continues

\item Nuno's complaints to Missoula PD met with intimidation (guns drawn)

\item Pattern: Institutional protection of YWCA-affiliated harassers

\end{itemize}

\section*{DETECTIVE BRUECKNER CONFLICT}

\begin{itemize}

\item \textbf{Role 1:} Missoula Police Detective (lead investigator)

\item \textbf{Role 2:} YWCA Missoula Board Member (simultaneous)

\item \textbf{Documented:} ProPublica records confirm board membership

\item \textbf{Disclosure:} NONE to defense counsel or court

\item \textbf{Legal Significance:} Unprecedented institutional capture

\end{itemize}

\section*{INSTITUTIONAL CORRUPTION}

\begin{itemize}

\item YWCA--Police nexus (board overlap)

\item Evidence suppression (character reference)

\item Officer recruitment (Ethan Smith)

\item Ongoing protection (2020--2025 harassment tolerance)

\end{itemize}

\section*{DAMAGES ATTRIBUTABLE TO E'LISE}

\begin{itemize}

\item Lost IT career: \$2,100,000

\item Specific contracts lost: \$680,000+

\item Legal fees: \$387,000

\item Medical/psychological: \$148,000

\item Business interference: \$500,000

\item Legal malpractice losses: \$6,400,000--\$8,400,000

\item \textbf{Total: \$10,215,000--\$12,215,000}

\end{itemize}

\section*{RECOMMENDED CHARGES}

\subsection*{Federal}

\begin{itemize}

\item 18 U.S.C. § 241: Conspiracy against rights (interstate with Danielle)

\item 18 U.S.C. § 242: Deprivation of rights under color of law (YWCA--police nexus)

\item 18 U.S.C. § 1341: Mail fraud (perjurious petition)

\item 18 U.S.C. § 1343: Wire fraud (false police reports)

\item 18 U.S.C. § 1503: Obstruction of justice (evidence suppression)

\item 18 U.S.C. § 1962: RICO predicate acts

\end{itemize}

\subsection*{State (Montana)}

\begin{itemize}

\item MCA 45-7-202: Perjury (written vs.\ sworn contradictions)

\item MCA 45-8-212: Criminal defamation (false criminal history claims)

\item MCA 45-4-102: Criminal conspiracy and solicitation (Tyleen Root)

\item MCA 45-5-203: Witness intimidation (continuing campaign)

\item MCA 45-5-204: False reporting (phantom calls)

\end{itemize}

\subsection*{Institutional Liability}

\begin{itemize}

\item YWCA Missoula (entity): Suppression of exculpatory evidence

\item City of Missoula (Monell liability): Detective Brueckner undisclosed conflict

\item Missoula Police Department: Pattern/practice institutional corruption

\end{itemize}

\section*{EVIDENCE AVAILABLE}

Court records, sworn testimony, written petition, phone records (no calls), character reference, ProPublica board documentation, judicial findings.

\end{document}
